%% ============================================================================
%% Section: Quantitative Layer (kept short: expressiveness, not a contribution)
%% Faithful to LinearGaussianInference.java, CounterfactualInference.java.
%% ============================================================================
\section{Quantitative Layer}\label{sec:quant}

The same compiled snapshot that answers structural queries also carries
optional numeric parameters, which let Causal4j answer quantitative causal
questions (interventions and counterfactuals) without leaving the DBMS. We
describe the layer briefly: it is a demonstration of the expressiveness the
indexed engine enables, not a contribution of the paper.

\noindent\textbf{Linear-Gaussian model.} When a model is parameterized, each
variable follows the structural equation
$V = \mu_{0,V} + \sum_{P \in \pa(V)} \beta_{P \rightarrow V}\,P + \varepsilon_V$
with $\varepsilon_V \sim \mathcal{N}(0, \sigma^2_V)$, where the edge
coefficients $\beta$, intercepts $\mu_0$, and noise variances $\sigma^2$ are
the \textsf{weight}, \textsf{intercept}, and \textsf{noise} properties read at
compile time (Section~\ref{sec:compile}). The parameters are supplied by the
upstream pipeline of Section~\ref{sec:lang:populate} or by the curator; the
DBMS never fits them. The snapshot reports itself parameterized only when some
value departs from the structural-only defaults, so a purely qualitative DAG
incurs no numeric machinery.

\noindent\textbf{Interventions and counterfactuals.} The do-operator
$do(X{=}x)$ fixes the mean of $X$ to $x$, sets its variance to zero, and drops
its incoming edges, which the general recursion then handles with no special
case. Interventional means propagate in topological order,
$\mu(V) = \mu_{0,V} + \sum_P \beta_{P \rightarrow V}\,\mu(P)$; interventional
variances follow from a covariance recursion in the same order; and the total
effect of $T$ on $Y$ is obtained by unit-impulse propagation
($\partial T = 1$, $\partial V = \sum_P \beta_{P \rightarrow V}\,\partial P$,
read off at $Y$), avoiding any path enumeration. Counterfactuals follow the
abduction, action, and prediction recipe over the same equations: infer the
exogenous noise consistent with the observed factual world, apply the
hypothesized intervention, and re-propagate. Every step reuses the
topological order already on the snapshot, so the quantitative layer adds
semantics, not a second data structure. On a parameterized version of the
running example,
\lstinline[style=causaldsl]|INTERVENE ON interventions = 1 COMPARE TO interventions = 0 RETURN EFFECT ON reported_cases|
returns the total effect along the directed paths, while a
\lstinline[style=causaldsl]|COUNTERFACTUAL| statement answers what
\textit{reported\_cases} would have been for one observed unit under a
different intervention. Since $do(\cdot)$ also yields a derived structural
world (Section~\ref{sec:incr:worlds}), the qualitative and quantitative
layers close the loop over a single stored graph.
